\documentclass[12pt,a4paper]{article}
\usepackage{fontspec}
\setmainfont{DejaVu Serif}
\setsansfont{DejaVu Sans}
\setmonofont{DejaVu Sans Mono}[Scale=0.85]
\usepackage[margin=2.2cm]{geometry}
\usepackage{amsmath}
\usepackage{xcolor}
\usepackage{enumitem}
\usepackage{booktabs}
\usepackage{array}
\usepackage{tikz}
\usepackage{hyperref}
\usepackage{titlesec}

\titleformat{\section}{\large\bfseries}{}{0pt}{}
\titleformat{\subsection}{\normalsize\bfseries}{}{0pt}{}
\setlength{\parskip}{0.5em}
\setlength{\parindent}{0pt}
\sloppy
\setlength{\emergencystretch}{3em}
\hypersetup{hidelinks}

\newcommand{\question}[2]{\section{Дасгал #1}\textbf{Асуулт.} #2\par\medskip\textbf{Хариулт.}\par}

% Gantt chart: \gantt{нэгжийн өргөн (cm)}{нэр/эхлэл/төгсгөл/өнгө, ...}{сүүлийн хугацаа}
\newcommand{\gantt}[3]{%
  \begin{center}
  \begin{tikzpicture}[x=#1cm, y=0.75cm]
    \foreach \n/\a/\b/\c in {#2} {
      \draw[fill=\c, draw=black!70] (\a,0) rectangle (\b,1);
      \node[font=\small] at ({(\a+\b)/2},0.5) {\n};
      \node[below, font=\scriptsize] at (\a,0) {\a};
    }
    \node[below, font=\scriptsize] at (#3,0) {#3};
  \end{tikzpicture}
  \end{center}}

\colorlet{cA}{blue!18}
\colorlet{cB}{green!20}
\colorlet{cC}{orange!25}
\colorlet{cD}{red!18}
\colorlet{cE}{violet!18}
\colorlet{cF}{teal!20}
\colorlet{cI}{gray!25}

\begin{document}

\begin{center}
  {\Large\bfseries Лабораторийн ажил №5}\\[0.4em]
  {\large CPU хуваарилалт — Бүлэг 5, Practice Exercises}\\[0.8em]
  F.CSM302 Үйлдлийн систем\\
  Гүйцэтгэсэн: Б.Билгүүн (B190910014)\\[0.3em]
  Сурах бичиг: A. Silberschatz, P. B. Galvin, G. Gagne, \textit{Operating System Concepts}, 10th ed., Wiley, 2018
\end{center}
\hrule
\bigskip

Тэмдэглэгээ: \emph{Эргэлтийн хугацаа} (turnaround time, TAT) $=$ дууссан хугацаа $-$ ирсэн хугацаа, \emph{Хүлээлтийн хугацаа} (waiting time, WT) $=$ TAT $-$ burst time.

%---------------------------------------------------------------
\question{5.1}{CPU хуваарилалтын алгоритм нь хуваарилагдсан процессуудын гүйцэтгэх дарааллыг тодорхойлдог. Нэг процессор дээр $n$ процессыг хуваарилах бол хэдэн өөр хуваарь (schedule) боломжтой вэ? $n$-ээр илэрхийлсэн томьёо гарга.}

$\mathbf{n!}$ ($n$ факториал) өөр хуваарь боломжтой.

Эхний байрлалд $n$ процессын аль нэгийг, хоёр дахь байрлалд үлдсэн $n-1$-ийн аль нэгийг гэх мэтээр сонгоно. Иймд нийт $n \cdot (n-1) \cdots 2 \cdot 1 = n!$ сэлгэмэл (permutation) гарна. Жишээ нь 5 процесст $5! = 120$ хуваарь боломжтой.

%---------------------------------------------------------------
\question{5.2}{Preemptive (буцаан авах) ба nonpreemptive (буцаан авахгүй) хуваарилалтын ялгааг тайлбарла.}

\begin{itemize}[leftmargin=1.5em]
  \item \textbf{Nonpreemptive} хуваарилалтад CPU авсан процесс түүнийг өөрөө чөлөөлөх хүртэл ашиглана. Процесс дуусах эсвэл I/O хүлээх (waiting төлөвт шилжих) үед л CPU чөлөөлөгдөнө. Хэрэгжүүлэхэд хялбар ч урт процесс богино процессуудыг удаан хүлээлгэж, хариу үйлдлийн хугацааг уртасгадаг.
  \item \textbf{Preemptive} хуваарилалтад үйлдлийн систем ажиллаж буй процессоос CPU-г хүчээр авч болно. Жишээ нь квант (time quantum) дуусах (timer тасалдал), эсвэл илүү өндөр эрэмбэтэй процесс ready болох үед ингэдэг. Интерактив болон бодит цагийн системд зайлшгүй хэрэгтэй. Харин хуваалцсан өгөгдөлд race condition үүсэх эрсдэлтэй тул синхрончлол шаарддаг.
\end{itemize}
Жишээ: FCFS нь nonpreemptive, RR нь preemptive. SJF болон priority хуваарилалт аль ч хувилбартай байж болно.

%---------------------------------------------------------------
\question{5.3}{Дараах процессууд заасан хугацаанд ирж, заасан хугацаагаар ажиллана. Nonpreemptive хуваарилалт ашиглаж, шийдвэр гаргах тухайн мөчид байгаа мэдээлэлд л тулгуурла.

\begin{center}
\begin{tabular}{ccc}\toprule
Процесс & Ирэх хугацаа & Burst time\\\midrule
$P_1$ & 0.0 & 8\\ $P_2$ & 0.4 & 4\\ $P_3$ & 1.0 & 1\\\bottomrule
\end{tabular}
\end{center}
(a) FCFS алгоритмаар дундаж эргэлтийн хугацаа хэд вэ? (b) SJF алгоритмаар хэд вэ? (c) SJF гүйцэтгэлийг сайжруулах ёстой ч бид $P_1$-ийг 0 мөчид эхлүүлсэн, учир нь удахгүй хоёр богино процесс ирэхийг мэдээгүй. CPU-г эхний 1 нэгж хугацаанд сул (idle) байлгаад дараа нь SJF хэрэглэвэл дундаж эргэлтийн хугацаа хэд болох вэ? Энэ хугацаанд $P_1$, $P_2$ хүлээж байгаа тул тэдний хүлээх хугацаа өсч болохыг санаарай. Энэ алгоритмыг future-knowledge scheduling гэж нэрлэж болно.}

\subsection{(a) FCFS}
\gantt{0.9}{$P_1$/0/8/cA, $P_2$/8/12/cB, $P_3$/12/13/cC}{13}
TAT: $P_1 = 8-0 = 8$,\; $P_2 = 12-0.4 = 11.6$,\; $P_3 = 13-1 = 12$.\\
Дундаж TAT $= (8 + 11.6 + 12)/3 = 31.6/3 \approx \mathbf{10.53}$.

\subsection{(b) SJF}
0 мөчид зөвхөн $P_1$ ирсэн тул түүнийг эхлүүлнэ. 8 мөчид $P_2$ (4), $P_3$ (1) хүлээж байгаа бөгөөд хамгийн богино нь $P_3$.
\gantt{0.9}{$P_1$/0/8/cA, $P_3$/8/9/cC, $P_2$/9/13/cB}{13}
TAT: $P_1 = 8$,\; $P_3 = 9-1 = 8$,\; $P_2 = 13-0.4 = 12.6$.\\
Дундаж TAT $= (8 + 8 + 12.6)/3 = 28.6/3 \approx \mathbf{9.53}$.

\subsection{(c) Эхний 1 нэгж idle, дараа нь SJF}
\gantt{0.9}{idle/0/1/cI, $P_3$/1/2/cC, $P_2$/2/6/cB, $P_1$/6/14/cA}{14}
TAT: $P_3 = 2-1 = 1$,\; $P_2 = 6-0.4 = 5.6$,\; $P_1 = 14-0 = 14$.\\
Дундаж TAT $= (1 + 5.6 + 14)/3 = 20.6/3 \approx \mathbf{6.87}$.

CPU 1 нэгж сул байсан ч богино процессуудыг түрүүлж ажиллуулснаар дундаж эргэлтийн хугацаа 9.53-аас 6.87 болж буурсан. Гэхдээ ирээдүйн ирэлтийг мэдэх боломжгүй тул бодит системд үүнийг хэрэгжүүлэх боломжгүй.

%---------------------------------------------------------------
\question{5.4}{Дараах процессуудын CPU burst time-ийг миллисекундээр өгөв. Процессууд $P_1, P_2, P_3, P_4, P_5$ дарааллаар бүгд 0 мөчид ирсэн гэж үзнэ.

\begin{center}
\begin{tabular}{ccc}\toprule
Процесс & Burst time & Priority\\\midrule
$P_1$ & 2 & 2\\ $P_2$ & 1 & 1\\ $P_3$ & 8 & 4\\ $P_4$ & 4 & 2\\ $P_5$ & 5 & 3\\\bottomrule
\end{tabular}
\end{center}
(a) FCFS, SJF, nonpreemptive priority (их тоо нь өндөр эрэмбэ) болон RR (квант $=2$) алгоритмуудаар ажиллуулахыг харуулсан дөрвөн Gantt chart зур. (b) Алгоритм тус бүрт процесс бүрийн эргэлтийн хугацаа хэд вэ? (c) Процесс бүрийн хүлээлтийн хугацаа хэд вэ? (d) Аль алгоритм хамгийн бага дундаж хүлээлтийн хугацаа өгөх вэ?}

\subsection{(a) Gantt chart}
\textbf{FCFS}
\gantt{0.75}{$P_1$/0/2/cA, $P_2$/2/3/cB, $P_3$/3/11/cC, $P_4$/11/15/cD, $P_5$/15/20/cE}{20}
\textbf{SJF}
\gantt{0.75}{$P_2$/0/1/cB, $P_1$/1/3/cA, $P_4$/3/7/cD, $P_5$/7/12/cE, $P_3$/12/20/cC}{20}
\textbf{Nonpreemptive priority} ($P_1$, $P_4$ хоёр ижил эрэмбэтэй тул FCFS дарааллаар)
\gantt{0.75}{$P_3$/0/8/cC, $P_5$/8/13/cE, $P_1$/13/15/cA, $P_4$/15/19/cD, $P_2$/19/20/cB}{20}
\textbf{RR (квант = 2)}
\gantt{0.75}{$P_1$/0/2/cA, $P_2$/2/3/cB, $P_3$/3/5/cC, $P_4$/5/7/cD, $P_5$/7/9/cE, $P_3$/9/11/cC, $P_4$/11/13/cD, $P_5$/13/15/cE, $P_3$/15/17/cC, $P_5$/17/18/cE, $P_3$/18/20/cC}{20}

\subsection{(b) Эргэлтийн хугацаа (мс)}
\begin{center}
\begin{tabular}{lccccc|c}\toprule
 & $P_1$ & $P_2$ & $P_3$ & $P_4$ & $P_5$ & Дундаж\\\midrule
FCFS     & 2  & 3  & 11 & 15 & 20 & 10.2\\
SJF      & 3  & 1  & 20 & 7  & 12 & 8.6\\
Priority & 15 & 20 & 8  & 19 & 13 & 15.0\\
RR       & 2  & 3  & 20 & 13 & 18 & 11.2\\\bottomrule
\end{tabular}
\end{center}

\subsection{(c) Хүлээлтийн хугацаа (мс)}
\begin{center}
\begin{tabular}{lccccc|c}\toprule
 & $P_1$ & $P_2$ & $P_3$ & $P_4$ & $P_5$ & Дундаж\\\midrule
FCFS     & 0  & 2  & 3  & 11 & 15 & 6.2\\
SJF      & 1  & 0  & 12 & 3  & 7  & \textbf{4.6}\\
Priority & 13 & 19 & 0  & 15 & 8  & 11.0\\
RR       & 0  & 2  & 12 & 9  & 13 & 7.2\\\bottomrule
\end{tabular}
\end{center}

\subsection{(d)}
\textbf{SJF} алгоритм хамгийн бага дундаж хүлээлтийн хугацаа (4.6 мс) өгнө. SJF нь дундаж хүлээлтийн хугацааг хамгийн бага болгодог оновчтой (optimal) алгоритм гэдэг нь онолоор батлагдсан.

%---------------------------------------------------------------
\question{5.5}{Дараах процессуудыг preemptive, round-robin хуваарилалтаар ажиллуулна.

\begin{center}
\begin{tabular}{cccc}\toprule
Процесс & Priority & Burst & Ирэх хугацаа\\\midrule
$P_1$ & 40 & 20 & 0\\ $P_2$ & 30 & 25 & 25\\ $P_3$ & 30 & 25 & 30\\
$P_4$ & 35 & 15 & 60\\ $P_5$ & 5 & 10 & 100\\ $P_6$ & 10 & 10 & 105\\\bottomrule
\end{tabular}
\end{center}
Процесс бүр тоон эрэмбэтэй бөгөөд их тоо нь өндөр эрэмбийг илэрхийлнэ. Эдгээрээс гадна систем CPU нөөц зарцуулдаггүй idle task ($P_{idle}$)-тай. Түүний эрэмбэ 0 бөгөөд өөр ажиллах процесс байхгүй үед хуваарилагдана. Квант (time quantum) 10 нэгж. Өндөр эрэмбэтэй процесс таслан зогсоовол зогсоогдсон процесс дарааллын төгсгөлд орно. (a) Gantt chart-аар хуваарилалтын дарааллыг харуул. (b) Процесс бүрийн эргэлтийн хугацаа хэд вэ? (c) Процесс бүрийн хүлээлтийн хугацаа хэд вэ? (d) CPU ашиглалтын хувь хэд вэ?}

\subsection{(a) Gantt chart}
\gantt{0.128}{$P_1$/0/20/cA, {\tiny idle}/20/25/cI, $P_2$/25/35/cB, $P_3$/35/45/cC, $P_2$/45/55/cB, $P_3$/55/60/cC, $P_4$/60/75/cD, $P_2$/75/80/cB, $P_3$/80/90/cC, idle/90/100/cI, $P_5$/100/105/cE, $P_6$/105/115/cF, $P_5$/115/120/cE}{120}

Алхам алхмаар:
\begin{itemize}[leftmargin=1.5em, itemsep=0pt]
  \item 0–20: $P_1$ ганцаараа тул квант дуусахад ч үргэлжлүүлэн ажиллаж, 20-д дуусна. 20–25: ажиллах процесс байхгүй тул $P_{idle}$ ажиллана.
  \item 25–35: $P_2$ ажиллана. 30-д ижил эрэмбэтэй $P_3$ ирж дараалалд орно. 35-д квант дуусахад $P_2$ (үлдэгдэл 15) дарааллын төгсгөлд орно.
  \item 35–45: $P_3$ (үлдэгдэл 15), 45–55: $P_2$ (үлдэгдэл 5), 55-аас $P_3$ ажиллана.
  \item 60: илүү өндөр эрэмбэтэй $P_4$ (35) ирж $P_3$-ийг таслан зогсооно. $P_3$-ийн үлдэгдэл 10 бөгөөд дарааллын төгсгөлд орно. 60–75: $P_4$ ажиллаж дуусна.
  \item 75–80: $P_2$ дуусна. 80–90: $P_3$ дуусна. 90–100: idle.
  \item 100–105: $P_5$ ажиллана. 105-д илүү өндөр эрэмбэтэй $P_6$ (10) ирж $P_5$-ийг таслан зогсооно. 105–115: $P_6$ дуусна. 115–120: $P_5$ дуусна.
\end{itemize}

\subsection{(b), (c) Эргэлтийн ба хүлээлтийн хугацаа}
\begin{center}
\begin{tabular}{lcccccc}\toprule
Процесс & Ирсэн & Burst & Дууссан & TAT & WT\\\midrule
$P_1$ & 0   & 20 & 20  & 20 & 0\\
$P_2$ & 25  & 25 & 80  & 55 & 30\\
$P_3$ & 30  & 25 & 90  & 60 & 35\\
$P_4$ & 60  & 15 & 75  & 15 & 0\\
$P_5$ & 100 & 10 & 120 & 20 & 10\\
$P_6$ & 105 & 10 & 115 & 10 & 0\\\midrule
Дундаж & & & & 30 & 12.5\\\bottomrule
\end{tabular}
\end{center}

\subsection{(d) CPU ашиглалт}
Нийт 120 нэгж хугацаанаас CPU 20–25 болон 90–100 хооронд, нийт 15 нэгж сул байсан. Ажилласан хугацаа нь бүх CPU ээлжийн (burst) нийлбэр $20+25+25+15+10+10 = 105$.
\[
  \text{CPU ашиглалт} = \frac{105}{120} = \mathbf{87.5\%}.
\]

%---------------------------------------------------------------
\question{5.6}{Олон түвшинт дараалал (multilevel queue) бүхий системийн өөр өөр түвшинд өөр хэмжээтэй квант (time quantum) ашиглах нь ямар давуу талтай вэ?}

Процессын төрөл бүрт тохирсон квант өгч болдог нь гол давуу тал юм.
\begin{itemize}[leftmargin=1.5em]
  \item \textbf{Өндөр эрэмбэтэй дараалал} (интерактив, I/O-bound процессууд) богино квант ашиглана. Эдгээр процесс CPU-г бага хугацаагаар хэрэглээд I/O руу шилждэг тул богино квант хариу өгөх хугацааг (response time) багасгана.
  \item \textbf{Бага эрэмбэтэй дараалал} (CPU-bound, batch процессууд) урт квант ашиглана. Эдгээр процесс CPU-г удаан ашигладаг тул урт квант нь орчин сэлгэлтийн (context switch) тоог, улмаар нэмэлт зардлыг (overhead) багасгаж, CPU-г үр ашигтай ашиглана.
\end{itemize}
Ингэснээр систем хариу үйлдлийн хурд болон үр ашгийг зэрэг хангаж чадна.

%---------------------------------------------------------------
\question{5.7}{CPU хуваарилалтын олон алгоритм параметртэй байдаг. Жишээ нь RR алгоритмд time slice-ийг заах параметр хэрэгтэй. Multilevel feedback queue-д дарааллын тоо, дараалал бүрийн хуваарилалтын алгоритм, процессыг дараалал хооронд шилжүүлэх шалгуур зэргийг тодорхойлох параметр хэрэгтэй. Тиймээс эдгээр нь үнэндээ алгоритмын олонлог юм (жишээ нь бүх time slice-тэй RR алгоритмуудын олонлог). Нэг олонлог нөгөөгөө агуулж болно: FCFS нь хязгааргүй time quantum-тэй RR алгоритм юм. Дараах хос алгоритмын олонлогуудын хооронд ямар хамаарал байх вэ (хэрэв байвал)? (a) Priority ба SJF, (b) Multilevel feedback queues ба FCFS, (c) Priority ба FCFS, (d) RR ба SJF}

\begin{enumerate}[label=(\alph*), leftmargin=1.5em]
  \item \textbf{Priority $\supset$ SJF.} SJF нь priority хуваарилалтын тусгай тохиолдол юм. Эрэмбийг дараагийн CPU ээлжийн (burst) таамагласан уртын урвуугаар тодорхойлдог: CPU ээлж богино байх тусам эрэмбэ өндөр.
  \item \textbf{MLFQ $\supset$ FCFS.} FCFS нь multilevel feedback queue-ийн тусгай тохиолдол юм. Ганцхан дараалалтай, хязгааргүй квант бүхий, процесс шилжүүлэлтгүй MLFQ нь FCFS-тэй яг адил ажиллана. Мөн MLFQ-ийн хамгийн доод дараалал ихэвчлэн FCFS байдаг.
  \item \textbf{Priority $\supset$ FCFS.} FCFS нь priority хуваарилалтын тусгай тохиолдол юм. Процесс эрт ирэх тусам эрэмбэ нь өндөр байхаар, эрэмбийг ирсэн хугацаагаар өгнө.
  \item \textbf{RR ба SJF: хамааралгүй.} Нэг нь нөгөөгөө агуулахгүй. RR нь CPU ээлжийн уртыг огт харгалздаггүй бол SJF нь зөвхөн CPU ээлжийн уртад тулгуурладаг. Аль ч параметртэй RR нь SJF-тэй ижил дараалал өгдөггүй, мөн эсрэгээрээ.
\end{enumerate}

%---------------------------------------------------------------
\question{5.8}{CPU хуваарилалтын алгоритм ойрын хугацаанд процессорыг хамгийн бага ашигласан процессуудыг илүүд үздэг гэж үзье. Энэ алгоритм яагаад I/O-bound програмуудыг илүүд үзэх боловч CPU-bound програмуудыг үүрд өлсгөхгүй (starvation) вэ?}

\begin{itemize}[leftmargin=1.5em]
  \item \textbf{I/O-bound програмыг илүүд үзэх шалтгаан.} I/O-bound процесс CPU-г богино ээлжээр (burst) хэрэглээд I/O хүлээх төлөвт шилждэг. Тиймээс ойрын хугацааны CPU ашиглалт нь бага байж, алгоритм түүнд өндөр эрэмбэ өгнө.
  \item \textbf{CPU-bound програм өлсөхгүй шалтгаан.}
  \begin{enumerate}[leftmargin=1.5em]
    \item I/O-bound процессууд ихэнх хугацааг I/O хүлээж өнгөрүүлдэг. Тэд блоклогдсон үед CPU сул болж, CPU-bound процесс ажиллах боломж олдоно.
    \item Алгоритм зөвхөн \emph{ойрын хугацааны} CPU ашиглалтыг тооцдог. CPU-bound процесс CPU авалгүй удаан хүлээх тусам түүний ойрын ашиглалтын утга буурч, эрэмбэ нь аажмаар өснө. Эцэст нь тэр CPU-г заавал авна. Энэ нь aging механизмтай ижил үр нөлөөтэй.
  \end{enumerate}
\end{itemize}

%---------------------------------------------------------------
\question{5.9}{PCS ба SCS хуваарилалтыг ялгаж тайлбарла.}

\begin{itemize}[leftmargin=1.5em]
  \item \textbf{PCS (process-contention scope).} Урсгалын сан хэрэглэгчийн түвшний урсгалуудыг сул LWP-үүд дээр хуваарилна. Энд CPU-ийн төлөөх өрсөлдөөн \emph{нэг процессын} урсгалуудын дунд явагдана. Many-to-one болон many-to-many загвартай системд хэрэглэгддэг бөгөөд ихэвчлэн программистын тогтоосон эрэмбээр хуваарилдаг.
  \item \textbf{SCS (system-contention scope, систем-өрсөлдөөний хүрээ).} Кернел кернелийн урсгалуудыг физик CPU дээр хуваарилна. Энд өрсөлдөөн \emph{систем дэх бүх} урсгалын дунд явагдана. One-to-one загвартай системүүд (Linux, Windows) зөвхөн SCS ашигладаг.
\end{itemize}
Pthreads-д \texttt{pthread\_attr\_setscope()} функцээр \texttt{PTHREAD\_SCOPE\_PROCESS} (PCS) эсвэл \texttt{PTHREAD\_SCOPE\_SYSTEM} (SCS)-ийг сонгоно.

%---------------------------------------------------------------
\question{5.10}{Уламжлалт UNIX хуваарилагч (scheduler) эрэмбийн тоо болон эрэмбэ хоёрын хооронд урвуу хамаарал тогтоодог: тоо их байх тусам эрэмбэ бага. Хуваарилагч секунд тутам процессуудын эрэмбийг дараах функцээр дахин тооцдог:
\[\text{Priority} = (\text{recent CPU usage} / 2) + \text{base},\]
энд base $= 60$. Recent CPU usage нь эрэмбийг өмнө нь тооцсоноос хойш процесс CPU-г хэр олон ашигласныг илэрхийлнэ. $P_1$-ийн recent CPU usage 40, $P_2$-ийнх 18, $P_3$-ийнх 10 гэж үзье. Эрэмбийг дахин тооцоход эдгээр гурван процессын шинэ эрэмбэ хэд болох вэ? Эндээс уламжлалт UNIX хуваарилагч CPU-bound процессын харьцангуй эрэмбийг өсгөдөг үү, бууруулдаг уу?}

\begin{center}
\begin{tabular}{lccc}\toprule
Процесс & Recent CPU usage & Тооцоо & Шинэ priority\\\midrule
$P_1$ & 40 & $40/2 + 60$ & \textbf{80}\\
$P_2$ & 18 & $18/2 + 60$ & \textbf{69}\\
$P_3$ & 10 & $10/2 + 60$ & \textbf{65}\\\bottomrule
\end{tabular}
\end{center}

Тоо их байх тусам эрэмбэ бага тул одоо $P_3$ хамгийн өндөр, $P_1$ хамгийн бага эрэмбэтэй болно.

Иймд уламжлалт UNIX хуваарилагч \textbf{CPU-bound процессын харьцангуй эрэмбийг бууруулдаг}. CPU-г их ашигласан процессын priority тоо өсч, эрэмбэ нь буурдаг. Харин CPU-г бага ашигладаг I/O-bound процессууд өндөр эрэмбэтэй болно.

\end{document}
